\chapter*{\IfLanguageName{dutch}{Woord vooraf}{Preface}}%
\label{ch:voorwoord}

Voor u ligt de bachelorproef die het sluitstuk vormt van mijn driejarige opleiding Toegepaste Informatica aan HOGENT. Dit onderzoek is het resultaat van een stageperiode bij het Italiaanse bedrijf Turtle Srl. De keuze voor dit specifieke onderwerp vloeide voort uit een concreet voorstel van mijn stagebedrijf. Bij de aanvang van dit traject bezat ik nog niet veel kennis over Artificiële Intelligentie -- en dan in het bijzonder over Retrieval-Augmented Generation (RAG) -- of over maatschappelijk verantwoord en duurzaam ondernemen. Juist vanwege die onbekendheid vond ik dit een mooie opportuniteit om mij te verdiepen in twee uiterst actuele en veelomvattende onderwerpen.

Met de uiteindelijke empirische resultaten uit deze scriptie krijgt Turtle Srl de gevalideerde, empirische resultaten in handen die gebruikt zullen worden om te beslissen welke specifieke chunkingstrategie, welk embeddingmodel en welk Large Language Model het beste geschikt zijn voor de realisatie van hun toekomstige Proof of Concept omtrent duurzaamheidsrapportage.

Deze stage bood mij tevens op persoonlijk vlak een unieke intellectuele uitdaging. Binnen mijn keuzetraject ligt de focus immers op Full Stack Development, waardoor dit project een uitdagende sprong in het AI-landschap betekende. Het gaf mij de kans om de theoretische en praktische werking van RAG-pipelines en evaluatiemethodieken volledig te doorgronden. Dit traject stelde mij in staat om de specifieke rollen en synergieën van verschillende architecturale componenten te ontleden, en bood diepgaand inzicht in hoe systeemconfiguraties de uiteindelijke prestaties en betrouwbaarheid van een volledige pipeline bepalen. Bovendien opende het onderzoek mijn ogen voor bedrijfsduurzaamheid: een domein waarmee je als IT-student tijdens de opleiding niet vaak in contact komt, maar dat desondanks een uiterst actueel en relevant onderwerp is waaraan in de toekomst een steeds groter strategisch belang zal worden gehecht binnen de convergentie met moderne AI-technologieën.

Het voltooien van deze bachelorproef was indirect een collectieve inspanning, en ik wil dan ook een aantal personen uitdrukkelijk bedanken voor hun onmisbare steun en begeleiding.

Allereerst gaat mijn oprechte dank uit naar mijn promotors vanuit HOGENT. Ik wil mevr. Nathalie Declercq hartelijk bedanken voor haar initiële begeleiding en de richting die zij aan de start van dit onderzoek heeft gegeven. Toen zij de opvolging van dit traject overdroeg, heeft mevr. Lena De Mol de rol van promotor op zich genomen. Haar feedback, methodologische sturing en continue ondersteuning tijdens het schrijfproces zijn van onschatbare waarde geweest. Daarnaast ben ik haar ontzettend dankbaar voor haar inspanningen om mij toegang te verlenen tot de Virtual IT Company (VIC) van HOGENT; dankzij haar snelle schakelen met de beheerders kreeg ik de kans om deze cloudinfrastructuur te gebruiken.

Daarnaast wil ik mijn co-promotor bij Turtle Srl, dhr. Gianluca Gueli, ten zeerste bedanken. Zijn begeleiding, technische expertise en heldere visie op de bedrijfsnoden hebben mij enorm geholpen om de theorie te vertalen naar concrete inzichten voor het bedrijf. In het bijzonder bood zijn eigen masterthesis over RAG-systemen — de allereerste bron die ik bij de aanvang van mijn onderzoek doornam — een fundamenteel referentiekader om kennis te maken met de materie en vormde de basis voor mijn verdere literatuurstudie. Tevens bedank ik het gehele Turtle-team en mijn tijdelijke collega's bij Turtle Srl voor de warme ontvangst en de fijne samenwerking binnen het bedrijf.

Een bijzonder woord van dank richt ik aan mijn medestudenten op buitenlandse stage: Alexandru Poenaru, Arnaud Tison, Tom Kluskens en Kyandro Voet. Onze brainstormsessies rond elkaars bachelorproeven waren niet alleen een bron van motivatie, maar hebben er ook concreet toe geleid dat technische problemen en bottlenecks tijdens het onderzoek efficiënt konden worden opgelost.

Ook wil ik de Virtual IT Company (VIC) bedanken voor het ter beschikking stellen van hun cloudinfrastructuur. Dankzij deze middelen was het mogelijk om het model \textit{gpt-oss-20b} binnen een private cloudomgeving te hosten, wat een waardevolle bijdrage leverde aan de evaluatiefase van dit onderzoek.

Tot slot wil ik mijn ouders en mijn zus uit de grond van mijn hart bedanken. Hun onvoorwaardelijke steun, geduld en voortdurende aanmoediging tijdens mijn gehele studententijd hebben mij de stabiele basis geboden om dit uitdagende doch prachtige studietraject met succes te voltooien. Een heel speciale vermelding gaat ook naar onze trouwe viervoeter, Stella; haar onuitputtelijke energie bracht de perfecte afleiding tijdens de meest stressvolle periodes, waardoor de academische druk even naar de achtergrond verdween.